\documentclass{article}
\usepackage{graphicx}
\usepackage{float} 
\usepackage{multirow}
\usepackage{pdfpages}
\usepackage{newtxtext,newtxmath}
\usepackage[a4paper, margin=1in]{geometry}
\usepackage[colorlinks=true,
            linkcolor=blue,
            citecolor=blue,
            urlcolor=blue]{hyperref}

\begin{document}

\hypersetup{pageanchor=false} 
\includepdf[
    pages=1,
    fitpaper=true,
    pagecommand={},
    linktodoc=true
]{Title page.pdf}
\clearpage
\hypersetup{pageanchor=true} 

\section*{Abstract}
\addcontentsline{toc}{section}{Abstract}

\large

Traditional sentiment analysis often classifies user reviews into simple, rigid categories such as positive, negative, or neutral. However, human emotions are inherently complex and nuanced. To better capture this depth, this thesis explores Dimensional Aspect-Based Sentiment Analysis (DimABSA) in the context of the SemEval-2026 Task 3. The primary objective is to extract specific product features (aspects), their broad categories, the expressed opinions, and precisely map their emotional intensity onto a continuous two-dimensional scale: Valence (positivity) and Arousal (intensity).
This thesis investigates all three subtasks of the SemEval-2026 DimABSA shared task, including dimensional sentiment regression, triplet extraction, and quadruplet extraction. Different transformer-based and Large Language Model (LLM)-based approaches were explored to address the challenges associated with continuous sentiment prediction and structured sentiment extraction.

For Subtask 1, multiple pretrained transformer encoder models, including BERT, RoBERTa, and DeBERTa variants, were evaluated for VA regression. Experimental results demonstrated that RoBERTa-large achieved the best performance across both restaurant and laptop domains, obtaining top leaderboard rankings in the official competition.

For Subtask 2, both an end-to-end Llama-based generative framework and a hybrid Llama + RoBERTa architecture were investigated for dimensional triplet extraction. Comparative experiments showed that the fully generative approach slightly outperformed the hybrid framework, highlighting the importance of preserving contextual dependencies during joint extraction and sentiment prediction.

For Subtask 3, instruction-tuned LLMs were employed for dimensional quadruplet extraction. Multiple experiments involving prompt engineering and hyperparameter optimization were conducted. The best-performing system achieved 7th rank on the official SemEval-2026 leaderboard for the laptop domain.

Overall, the findings demonstrate that transformer-based regression models and instruction-tuned LLMs provide effective solutions for Dimensional Aspect-Based Sentiment Analysis. The experiments further highlight the importance of prompt design, hyperparameter optimization, and parameter-efficient fine-tuning techniques for improving structured sentiment extraction performance.
\normalsize
\clearpage
\newpage
\tableofcontents
\newpage
\section{Introduction}
\label{sec:introduction}

In the era of e-commerce and digital communication, understanding user feedback is critical for product development and customer satisfaction. Traditional sentiment analysis typically classifies text at the document or sentence level into rigid, discrete categories such as \textit{positive}, \textit{negative}, or \textit{neutral}. However, human emotions are inherently complex and multifaceted. Aspect-Based Sentiment Analysis (ABSA) emerged to address this by identifying specific entities or features (aspects) mentioned in text and determining the sentiment directed toward each one \cite{pontiki2014semeval}. 

While traditional ABSA provides more granular insights, its reliance on categorical sentiment labels still limits its ability to capture the true intensity and nuance of human emotion. To overcome this limitation, recent research has shifted toward Dimensional Aspect-Based Sentiment Analysis (DimABSA). Grounded in psychology, particularly Russell's Circumplex Model of Affect \cite{russell1980circumplex}, dimensional sentiment maps emotions onto a continuous two-dimensional space: \textit{Valence} (ranging from extremely negative/unpleasant to extremely positive/pleasant) and \textit{Arousal} (ranging from calm/sleepy to excited/active). 

% NEW: Changed [htbp] to [H] to lock the image in place
\begin{figure}[H]
    \centering
    \includegraphics[width=0.8\textwidth]{VA.png} 
    \caption{Two-dimensional space of Valence and Arousal.}
    \label{fig:circumplex_model}
\end{figure}

This thesis focuses on advancing the state-of-the-art in DimABSA by participating in the SemEval 2026 DimABSA shared task (Track A: English). The ultimate objective of this research is to develop a highly accurate Information Extraction (IE) system capable of extracting full sentiment quadruplets from text: the Aspect Term, the Aspect Category, the Opinion Term, and the corresponding continuous Valence-Arousal (V-A) scores. To achieve this, the project leverages Large Language Models (LLMs)—specifically the Llama-3.1-8B architecture—optimised via Parameter-Efficient Fine-Tuning (PEFT) techniques such as Low-Rank Adaptation (LoRA).

This mid-thesis report details the progress made thus far and outlines the trajectory for the remainder of the research. At this current stage, the foundational data processing pipeline has been established, and the methodology for Subtask 1 (extracting Aspect Terms and V-A scores) has been successfully completed and evaluated. Building upon these promising preliminary results, the subsequent phases of this work will focus on tackling the more complex generative challenges of Subtask 2 (Category extraction) and the end-to-end quadruplet extraction required for Subtask 3.

\section{Problem Statement}
\label{sec:problem}

Dimensional Aspect-Based Sentiment Analysis (DimABSA) extends traditional Aspect-Based Sentiment Analysis (ABSA) by modelling sentiment as continuous values in the Valence–Arousal (VA) space rather than discrete polarity labels. The main goal of the SemEval 2026 DimABSA shared task is to teach computers how to truly understand the deep, complex emotions people share in customer reviews. Instead of just making a simple guess if a review is happy or sad, the system needs to find the exact features people are talking about and score their feelings using a continuous number scale for Valence (positivity) and Arousal (intensity). This work focuses on three subtasks defined in the SemEval-2026 shared task, each increasing in complexity from regression to structured information extraction.

\subsection{Subtask 1: Dimensional Aspect Sentiment Regression (DimASR)}

In this first step, the system is given a sentence and the specific aspect terms (the exact features) mentioned in it. The goal is for the model to figure out the exact Valence and Arousal (V-A) scores for each aspect. Unlike traditional sentiment classification, where polarity is assigned categorically, this task requires modelling fine-grained emotional intensity.

\textbf{Input:}
\begin{verbatim}
{
  "ID": "R001",
  "Text": "average to good thai food, but terrible delivery.",
  "Aspect": ["thai food", "delivery"]
}
\end{verbatim}

\textbf{Output:}
\begin{verbatim}
{
  "ID": "R001",
  "Aspect_VA": [
    {
      "Aspect": "thai food",
      "VA": "6.75#6.38"
    },
    {
      "Aspect": "delivery",
      "VA": "2.88#6.62"
    }
  ]
}
\end{verbatim}

\textbf{Challenges:}
\begin{itemize}
    \item Modelling sentiment as continuous values instead of discrete classes.
    \item Capturing context-dependent sentiment for multiple aspects within a single sentence.
    \item Learning the relationship between valence (polarity) and arousal (intensity).
    \item Handling domain-specific variations in expression (e.g., restaurant vs hotel reviews).
\end{itemize}

This subtask forms the core of the completed work in this thesis, where transformer-based regression models are employed to predict VA scores.

---

\subsection{Subtask 2: Dimensional Aspect Sentiment Triplet Extraction (DimASTE)}

Sometimes people do not mention a feature by name, but they talk about a general category instead. Subtask 2 extends the regression task by requiring the joint extraction of aspect terms, opinion terms, and their associated VA scores. The system must identify both the target (aspect) and the sentiment expression (opinion) directly from raw text.

\textbf{Input:}
\begin{verbatim}
{
  "ID": "R002",
  "Text": "The food was amazing but the service was slow."
}
\end{verbatim}

\textbf{Output:}
\begin{verbatim}
{
  "ID": "R002",
  "Triplet": [
    {
      "Aspect": "food",
      "Opinion": "amazing",
      "VA": "8.20#7.50"
    },
    {
      "Aspect": "service",
      "Opinion": "slow",
      "VA": "3.10#4.80"
    }
  ]
}
\end{verbatim}

\textbf{Challenges:}
\begin{itemize}
    \item Joint extraction of aspect and opinion spans from unstructured text.
    \item Ensuring correct alignment between aspects and their corresponding opinions.
    \item Generating structured outputs with strict formatting requirements.
    \item Simultaneously performing extraction and continuous regression.
\end{itemize}

This subtask is considered as future work in this thesis, where large language models are explored for structured generation.

---

\subsection{Subtask 3: Dimensional Aspect Sentiment Quadruplet Extraction (DimASQP)}

This is the final and hardest part of the project. The system is only given the raw text. It has to do all the heavy lifting on its own: find the aspect terms, map them to the right categories, pick out the exact opinion words, and finally, score the Valence and Arousal. 

\textbf{Input:}
\begin{verbatim}
{
  "ID": "R003",
  "Text": "The food quality was excellent but the service was disappointing."
}
\end{verbatim}

\textbf{Output:}
\begin{verbatim}
{
  "ID": "R003",
  "Quadruplet": [
    {
      "Aspect": "food quality",
      "Category": "FOOD#QUALITY",
      "Opinion": "excellent",
      "VA": "8.50#7.20"
    },
    {
      "Aspect": "service",
      "Category": "SERVICE#GENERAL",
      "Opinion": "disappointing",
      "VA": "2.70#5.60"
    }
  ]
}
\end{verbatim}

\textbf{Challenges:}
\begin{itemize}
    \item Incorporating domain-specific aspect categories with long-tail distributions.
    \item Increasing task complexity due to multi-component structured prediction.
    \item Maintaining consistency across extracted aspect, opinion, and category.
    \item Handling ambiguity in category assignment.
\end{itemize}

This subtask represents the most complex formulation and is planned as part of future work in this thesis.

---


\section{Literature Review}
\label{sec:literature}

\subsection{Aspect-Based Sentiment Analysis}

To truly understand how far Aspect-Based Sentiment Analysis (ABSA) has come, it helps to look back at how researchers first started teaching computers to understand human feelings. In the beginning, sentiment analysis treated whole reviews as just one feeling—either completely good or completely bad. However, as user reviews grew more detailed, researchers realised this was not enough. A person might love a laptop's screen but hate its battery. To solve this, the research community introduced Aspect-Based Sentiment Analysis (ABSA) during the early SemEval competitions between 2014 and 2016 \cite{pontiki2014semeval, pontiki2015semeval, pontiki2016semeval}, which introduced benchmark datasets for domains such as restaurant and laptop reviews. These tasks defined core components of ABSA, including aspect terms, opinion terms, and sentiment polarity.

Subsequent research extended ABSA from simple classification to structured prediction tasks. In particular, Aspect Sentiment Triplet Extraction (ASTE) and Aspect Sentiment Quadruplet Prediction (ASQP) were introduced to jointly extract aspect terms, opinion terms, and sentiment polarity, with ASQP further incorporating aspect category classification \cite{peng2020aste, cai2021asqp}. These early models were a huge step forward because they finally looked at specific features. Still, they forced human emotions into very rigid boxes: \textit{Positive}, \textit{Negative}, or \textit{Neutral}. While helpful, these simple labels missed the true depth of how people actually feel.

As noted in the SemEval-2026 task description, ABSA has evolved significantly, with applications spanning multiple domains such as finance, healthcare, and social media \cite{yu2026dimabsa}. However, the reliance on categorical sentiment labels restricts the ability of these systems to capture nuanced emotional expressions.

\subsection{Dimensional Sentiment Analysis}

Human emotions are rarely just black and white. Sometimes we are intensely angry, and other times we are just mildly annoyed. To overcome the limitations of categorical sentiment representation, dimensional sentiment analysis has been proposed as an alternative framework. Grounded in psychological theories such as Russell’s Circumplex Model of Affect, emotions are represented along continuous dimensions, primarily valence and arousal \cite{russell1980circumplex}. This representation allows for capturing both the polarity and intensity of sentiment.

Several resources and datasets have been developed to support dimensional sentiment analysis. Lexicon-based approaches such as NRC-VAD provide word-level valence, arousal, and dominance scores \cite{mohammad2018nrcvad}, while sentence-level datasets such as EmoBank extend this to larger textual units \cite{buechel2017emobank}. As highlighted in the SemEval-2026 organisers' paper, dimensional representations enable more fine-grained distinctions in emotional expression and support broader applications beyond traditional sentiment classification \cite{yu2026dimabsa}.

The integration of dimensional sentiment into ABSA has led to the formulation of Dimensional Aspect-Based Sentiment Analysis (DimABSA), which combines structured extraction with continuous regression. This paradigm shift represents a significant advancement in sentiment analysis research.

\subsection{Transformer-Based Approaches}

Recent advances in NLP have been driven by pretrained transformer models such as BERT \cite{devlin2019bert}, RoBERTa \cite{liu2019roberta}, and DeBERTa \cite{he2021deberta}. These models have demonstrated strong performance across a wide range of tasks, including sentiment analysis and ABSA.

In the context of DimABSA, transformer-based encoders are commonly used for regression tasks. As reported in the SemEval-2026 shared task, many top-performing systems rely on fine-tuned transformer architectures to predict valence and arousal scores \cite{yu2026dimabsa}. Ensemble methods, combining multiple encoder models, further improve performance by reducing variance and increasing robustness.

Our approach for Subtask 1 follows this paradigm by employing a weighted ensemble of transformer-based models. Similar strategies have been shown to enhance predictive accuracy, particularly in multi-target regression settings, where both valence and arousal must be estimated simultaneously.

\subsection{Large Language Models for Structured Extraction}

While encoder-based models are effective for regression tasks, structured extraction tasks such as triplet and quadruplet prediction require more flexible modelling capabilities. Recent research has explored the use of instruction-tuned Large Language Models (LLMs), such as Llama \cite{touvron2023llama}, for these tasks.

LLMs enable end-to-end generation of structured outputs, allowing models to simultaneously extract aspect terms, opinion terms, and sentiment scores. As discussed in recent system approaches, generative models can overcome limitations of pipeline-based methods, which often suffer from error propagation across stages.

In particular, fine-tuning LLMs for structured prediction allows the model to learn the relationship between textual expressions and their corresponding sentiment intensities in a unified framework. This approach is especially effective for complex tasks such as quadruplet extraction, where both categorical and continuous outputs must be generated.

\subsection{Challenges in DimABSA}

Despite recent progress, DimABSA remains a challenging problem. The SemEval-2026 organisers highlight several key difficulties, including the need to jointly model structured extraction and continuous regression, as well as the complexity introduced by multilingual and multi-domain datasets \cite{yu2026dimabsa}.

Additionally, evaluation of such systems requires specialised metrics. The proposed continuous F1 (cF1) metric integrates both categorical correctness and regression accuracy, reflecting the hybrid nature of the task. However, optimising for such metrics remains non-trivial.

Another challenge lies in capturing the relationship between valence and arousal, which are interdependent dimensions. Models must not only predict accurate values but also maintain consistency within the emotional space.

\section{Datasets}
\label{sec:datasets}

This work utilises the datasets provided as part of the SemEval-2026 Task 3 on Dimensional Aspect-Based Sentiment Analysis (DimABSA). The datasets are designed to support both regression and structured extraction tasks across multiple languages and domains. In this thesis, experiments are conducted on the \textbf{English language datasets}, specifically focusing on the \textbf{restaurant} and \textbf{laptop} domains.

\subsection{Dataset Overview}

The DimABSA dataset is a multilingual and multi-domain benchmark that extends traditional ABSA annotations by incorporating continuous valence–arousal (VA) scores. Each instance is annotated with different levels of granularity depending on the subtask, including aspect terms, opinion terms, aspect categories, and VA scores \cite{lee2026dimabsabuildingmultilingualmultidomain}.

As reported in the organisers' paper, Track A datasets span multiple domains such as restaurant, laptop, hotel, and finance, and include annotations for up to four components: Aspect (A), Category (C), Opinion (O), and Valence–Arousal (VA). The VA scores are continuous values ranging from 1 to 9, where 1 represents extreme negativity or low arousal, and 9 represents extreme positivity or high arousal.

\subsection{Domains Used in This Work}

This thesis focuses on two primary domains:

\textbf{Restaurant Domain:}  
The restaurant dataset contains user reviews describing food quality, service, ambience, and other related aspects. This domain is widely used in ABSA research due to its rich and diverse sentiment expressions. It includes multiple aspects per sentence, often with contrasting sentiments (e.g., positive food quality but negative service).

\textbf{Laptop Domain:}  
The laptop dataset consists of product reviews focusing on features such as battery life, screen quality, performance, and price. Compared to the restaurant domain, this dataset often contains more technical and feature-specific language, making it challenging for models to capture sentiment nuances.

\subsection{Data Characteristics}

The datasets exhibit several important characteristics:

\begin{itemize}
    \item \textbf{Multi-aspect Sentences:} A single sentence may contain multiple aspects with different sentiment expressions.
    \item \textbf{Continuous Sentiment Representation:} Sentiment is represented using continuous VA scores instead of discrete labels.
    \item \textbf{Domain Diversity:} Different domains exhibit distinct linguistic patterns and sentiment distributions.
    \item \textbf{Structured Annotations:} Higher-level subtasks require joint modelling of multiple elements, increasing complexity.
\end{itemize}

According to the organisers' paper, the dataset contains tens of thousands of annotated aspect instances across multiple languages and domains. Together with parallel initiatives like the DimStance dataset, this makes it one of the most comprehensive benchmarks for dimensional sentiment analysis \cite{lee2026dimabsabuildingmultilingualmultidomain}.

\subsection{Annotation Process}

The dataset annotation follows a two-stage process. First, categorical elements such as aspect terms, opinion terms, and aspect categories are identified. Subsequently, valence and arousal scores are assigned by multiple annotators, and the final score is obtained by averaging their ratings. This process ensures both structural consistency and reliable sentiment estimation \cite{lee2026dimabsabuildingmultilingualmultidomain}.

\section{Methodology}
\label{sec:methodology}

This section presents the methodologies adopted for all three subtasks of the SemEval-2026 Dimensional Aspect-Based Sentiment Analysis (DimABSA) shared task. The proposed approaches combine transformer-based regression models and instruction-tuned Large Language Models (LLMs) to address both continuous sentiment prediction and structured sentiment extraction. Multiple pretrained transformer architectures, including BERT, RoBERTa, and DeBERTa variants, were explored for dimensional sentiment regression, while instruction-tuned LLMs such as Llama, Qwen, and DeepSeek-R1 are utilized for triplet and quadruplet extraction tasks. The following subsections describe the methodologies, model architectures, and training strategies employed for each subtask.

\subsection{Subtask 1: Dimensional Aspect Sentiment Regression}

Subtask 1 is formulated as a multi-target regression problem, where the objective is to predict continuous valence and arousal scores for given aspect terms within a sentence. To solve this task, transformer-based encoder models are employed due to their strong contextual representation capabilities.

\subsubsection{Model Architecture}

The regression model is based on pretrained transformer encoders. Given an input sentence and an associated aspect term, the input is formatted as a combined sequence and passed through the encoder. The contextual representation of the sequence is then used to predict valence and arousal scores.

Formally, given an input sequence $x$, the encoder produces a contextual representation $h$. A regression head is applied to obtain the final predictions:

\begin{equation}
\hat{y} = W h + b
\end{equation}

where $\hat{y} = (V, A)$ represents the predicted valence and arousal scores, and $W, b$ are learnable parameters.

\subsubsection{Model Exploration and Selection}

Multiple pretrained transformer models were explored during experimentation, including BERT, RoBERTa, and DeBERTa variants. These models were fine-tuned on the task-specific dataset to evaluate their performance on valence-arousal prediction.

Empirical results indicated that \textbf{RoBERTa-large} consistently outperformed other models across both restaurant and laptop domains. This improvement is attributed to its robust pretraining strategy and ability to capture nuanced contextual information. Consequently, RoBERTa-large was selected as the primary model for the final system.

\subsubsection{Training Strategy}

The models were fine-tuned using supervised learning on the provided training data. The training process involves:

\begin{itemize}
    \item Tokenization of input text using the corresponding tokenizer.
    \item Encoding of aspect-aware input sequences.
    \item Optimization using a regression loss function (e.g., Mean Squared Error).
    \item Validation on a development set to monitor performance.
\end{itemize}

The model is trained to jointly predict both valence and arousal values, enabling it to learn the relationship between sentiment polarity and intensity.

\subsubsection{Evaluation Metric}

Subtask 1 (DimASR) is formulated as a sentiment regression task and is evaluated using the \textbf{Root Mean Square Error (RMSE)} between the predicted and gold valence–arousal (VA) scores. The official evaluation metric used in the SemEval-2026 shared task is defined as follows:

\begin{equation}
RMSE_{VA} =
\sqrt{
\frac{
\sum_{i=1}^{N}
\left[
\left(V_p^{(i)} - V_g^{(i)}\right)^2
+
\left(A_p^{(i)} - A_g^{(i)}\right)^2
\right]
}{N}
}
\end{equation}

where:
\begin{itemize}
    \item $N$ denotes the total number of instances,
    \item $V_p^{(i)}$ and $A_p^{(i)}$ represent the predicted valence and arousal scores for the $i^{th}$ instance,
    \item $V_g^{(i)}$ and $A_g^{(i)}$ represent the corresponding ground-truth valence and arousal scores.
\end{itemize}

The RMSE metric measures the average prediction error between the predicted and gold VA values. Lower RMSE values indicate better model performance, as they correspond to smaller deviations from the ground-truth emotional scores.

Additionally, the predicted VA outputs are constrained within the range $[1, 9]$ and are rounded to two decimal places, following the official task requirements.

\subsection{Subtask 2: Dimensional Aspect Sentiment Triplet Extraction (DimASTE)}
\label{subsec:subtask2_method}

Subtask 2 was formulated as a structured information extraction problem, where the objective was to extract sentiment triplets consisting of \textit{(Aspect, Opinion, VA)} directly from raw text. To address this task, two different experimental approaches were explored: (1) an end-to-end LLM-based triplet generation approach, and (2) a hybrid approach combining LLM-based extraction with transformer-based VA regression.

\subsubsection{Dataset Preparation and Prompt Formulation}

Experiments were conducted on the English restaurant and laptop datasets provided in the SemEval-2026 DimABSA shared task. The datasets were loaded in JSON Lines format and processed separately for each domain. Each sample contained the review text and corresponding triplet annotations consisting of aspect terms, opinion terms, and valence–arousal scores.

To prepare the data for instruction tuning, each training instance was converted into a strict instruction-response format. This formulation enabled the model to learn structured generation in a zero-shot or few-shot inference setting. A unified prompt template was engineered to enforce strict extraction rules and formatting constraints. The exact prompt template utilized for this task was defined as follows:

\begin{verbatim}
Below is an instruction describing a task, paired with an input that provides 
additional context. Your goal is to generate an output that correctly 
completes the task.

### Instruction:
Given a textual instance [Text], extract all (A, O, VA) triplets, where:
- A is an Aspect term (the entity being discussed)
- O is an Opinion term (the feeling or sentiment expressed about A)
- VA is a Valence-Arousal score in the format (valence#arousal)

Valence: 1.00 (negative) to 9.00 (positive).
Arousal: 1.00 (calm) to 9.00 (excited).
Important: Aspect and Opinion must maintain the exact case and spelling 
found in the [Text].

### Example:
Input:
[Text] average to good thai food, but terrible delivery.

Output:
[Triplet] (thai food, average to good, 6.75#6.38), (delivery, terrible, 2.88#6.62)

### Question:
Now complete the following example:
Input:
{Review Text Inserted Here}
\end{verbatim}

\subsubsection{Experiment 1: End-to-End LLM-Based Triplet Extraction}

In the first experiment, the entire triplet generation process was performed using the \textbf{Llama-3.1-8B-Instruct} model. The model was fine-tuned using the Unsloth framework with 4-bit quantization to reduce GPU memory requirements and enable efficient training.

Parameter-Efficient Fine-Tuning (PEFT) was applied using the LoRA technique. LoRA adapters were injected into the attention projection layers, including the query, key, value, and output projection modules. 

\paragraph{Training Loop Architecture}
The training process was orchestrated using the \texttt{SFTTrainer} (Supervised Fine-Tuning Trainer) from the Hugging Face \texttt{trl} library. Because this was formulated as a causal language modeling task, the model was optimized using standard next-token prediction cross-entropy loss, strictly calculated only on the generated triplet output (the prompt tokens were masked during loss calculation). 

To ensure stability while training an 8-billion parameter model on constrained hardware, the training loop utilized gradient checkpointing and mixed precision (\texttt{fp16}/\texttt{bf16}). Optimization was handled by an 8-bit AdamW optimizer (\texttt{adamw\_8bit}) paired with a linear learning rate scheduler. To prevent sequence length mismatch errors during batching, all input sequences were padded or truncated to a fixed maximum sequence length. 

The model was trained separately for the laptop and restaurant domains. During inference, the model generated structured triplets directly from the input text. A regex-based post-processing step was then applied to extract valid triplets in the final target format:
\begin{equation}
(\text{Aspect}, \text{Opinion}, \text{VA})
\end{equation}
where VA is represented using the \texttt{V\#A} format. This approach treated the task as a fully generative extraction problem, allowing the model to jointly learn aspect extraction, opinion extraction, and continuous sentiment prediction within a single pass.

\subsubsection{Experiment 2: Hybrid (LLM + RoBERTa) Approach}

This subtask involves two major components: textual extraction and VA score generation, where VA prediction is fundamentally a regression task. During experimentation, it was observed that although Llama-based models perform well in extracting textual information such as aspect and opinion terms, they are comparatively less reliable for precise continuous regression tasks, particularly when generating decimal-valued VA scores. In contrast, RoBERTa models are widely recognised for their strong performance in regression-based tasks. Therefore, a hybrid architecture was explored to investigate whether combining the strong extraction capabilities of Llama with the robust regression performance of RoBERTa could improve the overall triplet extraction and VA prediction performance.

In this approach, the \textbf{Llama-3.1-8B-Instruct} model was utilized strictly for extracting the textual aspect and opinion terms. The continuous numerical VA scores were then predicted separately using a fine-tuned \textbf{RoBERTa-large} regression model.

First, Llama generated structured aspect–opinion pairs from the input text using a simplified version of the unified prompt. These extracted terms were then concatenated into a cross-encoder format and passed to RoBERTa-large:
\begin{equation}
\text{Text} \; \text{\texttt{</s>}} \; \text{Aspect} \; \text{\texttt{</s>}} \; \text{Opinion}
\end{equation}

The RoBERTa-large model was trained as a multi-target regression model to jointly predict valence and arousal scores, outputting two continuous floating-point values corresponding to the exact VA dimensions.

This hybrid framework effectively combined the strengths of both architectures:
\begin{itemize}
    \item \textbf{Llama-3.1-8B} provided strong generative and structured natural language extraction capabilities.
    \item \textbf{RoBERTa-large} provided highly stable and mathematically accurate continuous regression performance.
\end{itemize}

\subsection{Subtask 3: Dimensional Aspect Sentiment Quadruplet Extraction (DimASQP)}

Subtask 3 extends the DimABSA framework by requiring the extraction of complete sentiment quadruplets consisting of \textit{(Aspect, Category, Opinion, VA)} directly from raw text. Compared to the previous subtasks, this task is significantly more challenging because it combines structured information extraction, category classification, and continuous valence–arousal prediction within a unified framework.

All experiments for this subtask were conducted exclusively on the \textbf{English laptop domain} dataset provided in the SemEval-2026 shared task.

\subsubsection{Dataset Preparation and Data Extraction}

The dataset was provided in JSON Lines format, where each instance contains the review text and corresponding quadruplet annotations. Each quadruplet consists of an aspect term, an aspect category, an opinion term, and a valence–arousal score represented in the \texttt{V\#A} format.

To prepare the dataset for instruction tuning, the annotations were converted into structured target sequences. Each input sample was transformed into an instruction-response format, where the instruction contains the review text and the expected output contains the list of quadruplets.

The preprocessing pipeline involved:
\begin{itemize}
    \item Loading and parsing JSONL files.
    \item Extracting aspect, category, opinion, and VA annotations.
    \item Formatting annotations into structured textual outputs.
    \item Creating instruction-based prompts for supervised fine-tuning.
\end{itemize}

This formulation enables the task to be treated as a sequence generation problem using instruction-tuned Large Language Models.

\subsubsection{Model Architecture}

The experiments were conducted using the \textbf{Llama-3.1-8B-Instruct} model. To enable efficient training within limited computational resources, the model was loaded using 4-bit quantization through the Unsloth framework.

Parameter-Efficient Fine-Tuning (PEFT) was applied using the LoRA (Low-Rank Adaptation) technique. LoRA adapters were inserted into the transformer attention projection layers, including:
\begin{itemize}
    \item Query projection
    \item Key projection
    \item Value projection
    \item Output projection
\end{itemize}

The training process utilized:
\begin{itemize}
    \item Mixed precision training
    \item Gradient checkpointing
    \item Supervised fine-tuning
    \item Instruction-based prompting
\end{itemize}

The model was trained to directly generate structured quadruplets from raw review text in an end-to-end manner.

\subsubsection{Training Loop}

The training loop follows a supervised instruction-tuning framework. During training, each input review is tokenized and passed to the Llama model together with the target quadruplet sequence.

The workflow consists of the following steps:
\begin{enumerate}
    \item Tokenization of input prompts and target outputs.
    \item Generation of input IDs and attention masks.
    \item Forward propagation through the Llama model.
    \item Computation of sequence generation loss.
    \item Backpropagation and parameter updates using AdamW optimization.
\end{enumerate}

The model was trained for multiple epochs with periodic evaluation and checkpoint saving. During inference, generated outputs were post-processed using regex-based extraction to obtain valid quadruplets.

\subsubsection{Experiment 1: Baseline Llama Quadruplet Extraction}

The first experiment focused on directly generating quadruplets using a highly detailed, descriptive instruction prompt. In this setup, the Llama-3.1-8B-Instruct model was fine-tuned to learn the complete mapping from the raw review text to structured quadruplet outputs.

This experiment established the baseline performance for the task and evaluated the ability of instruction-tuned LLMs to jointly model textual extraction and numerical regression. The baseline prompt utilized a few-shot approach, explicitly detailing step-by-step reasoning and providing concrete examples to guide the model:

\begin{verbatim}
You are an expert Linguist specializing in Aspect-Based Sentiment Analysis 
(ABSA). Your task is to extract highly accurate (A, C, O, VA) quadruplets 
from the given text.

### Core Extraction Rules:
1. Aspect (A): The specific feature or entity mentioned. Must match the 
   input text casing exactly.
2. Category (C): Classify the aspect using the "ENTITY#ATTRIBUTE" schema 
   below. Use ONLY these labels. Must be UPPERCASE.
3. Opinion (O): The specific word/phrase used to express the sentiment. 
   Match the input casing exactly.
4. Valence-Arousal (VA): 
   - Valence: 1.00 (Extremely Negative) to 9.00 (Extremely Positive). 
     5.00 is Neutral.
   - Arousal: 1.00 (Calm/Sleepy) to 9.00 (Excited/Angry). 5.00 is Moderate.
   - Format as "V.VV#A.AA" (always 2 decimal places).

### Label Schema Constraints:
- Valid Entities: {laptop_entity}
- Valid Attributes: {laptop_attribute}

### Step-by-Step Reasoning:
Step 1: Identify all sentiment-bearing phrases and the aspects they refer to.
Step 2: Map each aspect to the most relevant ENTITY and ATTRIBUTE.
Step 3: Determine the numerical Valence and Arousal of the opinion.
Step 4: Format as a list of (A, C, O, VA) quadruplets.

### Examples:
Input: [Text] The screen is incredibly bright and vibrant, but the price 
is a bit steep.
Output: [Quadruplet] (screen, DISPLAY#QUALITY, incredibly bright and vibrant, 
8.50#7.20), (price, LAPTOP#PRICE, a bit steep, 3.20#5.50)

### Target Task:
Input:
[Text] {input_text}
Output:
\end{verbatim}

\subsubsection{Experiment 2: Prompt Change}

While detailed prompts are generally helpful for LLMs, excessively long instructions can sometimes dilute the model's focus during fine-tuning, leading to hallucinated formatting or skipped categories. In the second experiment, the prompt structure was significantly condensed to improve output consistency and extraction quality.

The modified prompt removed the conversational framing, the step-by-step reasoning, and the few-shot examples. Instead, it relied on a strict, minimalist set of formatting instructions. This zero-shot style prompt was designed to force the model to focus purely on the target rules:

\begin{verbatim}
Please analyze the laptop review below and extract all sentiment quadruplets 
in the format (Aspect, Category, Opinion, Valence#Arousal).

Rules:
- Aspect: The exact words from the text representing the feature.
- Category: Must be chosen from the allowed list (ENTITY#ATTRIBUTE).
- Opinion: The exact words expressing the sentiment.
- Valence#Arousal: Numbers from 1.00 to 9.00 formatted as V.VV#A.AA.

Allowed Categories:
Entities: {laptop_entity}
Attributes: {laptop_attribute}

Review:
[Text] {input_text}

Output:
\end{verbatim}

This experiment aimed to study the direct impact of prompt engineering on structured sentiment extraction performance, specifically testing whether a strict, concise instruction would yield fewer hallucinations and higher format adherence than a verbose, step-by-step instruction.

\subsubsection{Experiment 3: Hyperparameter Optimization}

The third experiment focused on optimizing training hyperparameters to further improve generation quality and model stability. Several hyperparameters were modified during experimentation, including:
\begin{itemize}
    \item Learning rate
    \item Batch size
    \item LoRA rank
    \item Number of training epochs
    \item Maximum sequence length
\end{itemize}

The objective of this experiment was to analyze how training configurations influence quadruplet extraction performance and structured output consistency.

\subsubsection{Evaluation Metrics for Subtask 2 and 3}

Subtask 2 (DimASTE) and Subtask 3 (DimASQP) are hybrid sentiment analysis tasks that combine both structured extraction and continuous regression. The outputs contain categorical components such as Aspect (A), Category (C), and Opinion (O), together with continuous Valence–Arousal (VA) scores. Traditional ABSA evaluation metrics typically assess categorical elements using Precision, Recall, and F1-score, where a prediction is considered correct only if all categorical elements exactly match the gold annotation. However, such binary evaluation does not account for the quality of continuous VA predictions.

To address this limitation, the SemEval-2026 shared task introduces the \textbf{continuous True Positive (cTP)} metric, which extends the traditional categorical true positive by incorporating a penalty based on the VA prediction error.

Let $P$ denote the set of predicted triplets or quadruplets. For a prediction $t \in P$, the continuous true positive is defined as:

\begin{equation}
cTP(t)=
\begin{cases}
1 - dist(VA_p(t), VA_g(t)), & t \in P_{cat} \\
0, & \text{otherwise}
\end{cases}
\end{equation}

where:
\begin{itemize}
    \item $P_{cat} \subseteq P$ represents the set of predictions whose categorical components exactly match the gold annotations,
    \item $VA_p(t)$ denotes the predicted valence–arousal pair,
    \item $VA_g(t)$ denotes the corresponding gold valence–arousal pair.
\end{itemize}

The distance between predicted and gold VA values is computed using the normalized Euclidean distance:

\begin{equation}
dist(VA_p, VA_g)=
\frac{
(V_p - V_g)^2 + (A_p - A_g)^2
}{D_{max}}
\end{equation}

where:
\begin{itemize}
    \item $V_p$ and $A_p$ represent predicted valence and arousal scores,
    \item $V_g$ and $A_g$ represent gold valence and arousal scores,
    \item $D_{max} = 8^2 + 8^2 = 128$ is the maximum possible squared Euclidean distance in the VA space on the $[1,9]$ scale.
\end{itemize}

This normalization ensures that:

\begin{equation}
dist(\cdot) \in [0,1]
\end{equation}

Building upon the continuous true positive, the \textbf{continuous Recall (cRecall)} is defined as:

\begin{equation}
cRecall =
\frac{
TP_{cat} -
\sum_{t \in P_{cat}}
dist(VA_p(t), VA_g(t))
}{
TP_{cat} + FN_{cat}
}
\end{equation}

where:
\begin{itemize}
    \item $TP_{cat}$ denotes the number of categorically correct predictions,
    \item $FN_{cat}$ denotes the number of gold annotations with no categorical match.
\end{itemize}

Similarly, the \textbf{continuous Precision (cPrecision)} is computed as:

\begin{equation}
cPrecision =
\frac{
TP_{cat} -
\sum_{t \in P_{cat}}
dist(VA_p(t), VA_g(t))
}{
TP_{cat} + FP_{cat}
}
\end{equation}

where $FP_{cat}$ represents the number of predicted tuples with no categorical match.

Finally, the official evaluation metric for Subtasks 2 and 3 is the \textbf{continuous F1-score (cF1)}, defined as the harmonic mean of cRecall and cPrecision:

\begin{equation}
cF1 =
\frac{
2 \times cRecall \times cPrecision
}{
cRecall + cPrecision
}
\end{equation}

This metric provides a unified evaluation framework that jointly measures:
\begin{itemize}
    \item correctness of categorical extraction,
    \item quality of aspect and opinion identification,
    \item category classification accuracy,
    \item precision of continuous valence–arousal prediction.
\end{itemize}

Compared to traditional binary F1 metrics, cF1 offers a more realistic assessment of DimABSA systems by considering both structured extraction quality and continuous sentiment prediction performance simultaneously.

\section{Results}
\label{sec:results}

\subsection{Subtask 1 Result}

To evaluate the performance of different transformer-based models for Subtask 1 (Dimensional Aspect Sentiment Regression), multiple pretrained architectures were fine-tuned and compared using Mean Squared Error (MSE) on the validation set.

\begin{table}[htbp]
\centering
\caption{Performance comparison of different models on the validation set (lower MSE indicates better performance).}
\label{tab:model_leaderboard}
\begin{tabular}{lc}
\hline
\textbf{Model} & \textbf{Best Validation Loss (MSE)} \\
\hline
\textbf{RoBERTa-large} & \textbf{0.7181} \\
BERT-base-uncased & 0.9016 \\
BERT-base-cased & 0.9147 \\
RoBERTa-base & 1.0214 \\
DeBERTa-v3-base & 1.0359 \\
\hline
\end{tabular}
\end{table}

The results clearly show that \textbf{RoBERTa-large achieves the lowest validation loss}, significantly outperforming other models. This demonstrates its superior ability to capture contextual information and model continuous sentiment representations effectively.

\subsection{Official Evaluation Results}

The final RoBERTa-large model was used to generate predictions, which were submitted to the official SemEval-2026 evaluation platform. The performance was evaluated using RMSE on the valence-arousal predictions.

\begin{table}[htbp]
\centering
\caption{Official evaluation results for Subtask 1 (DimASR). Lower RMSE indicates better performance.}
\label{tab:official_results}
\begin{tabular}{lccc}
\hline
\textbf{Domain} & \textbf{RMSE\_VA} & \textbf{PCC\_V} & \textbf{PCC\_A} \\
\hline
Laptop (ENG) & 1.3283 & 0.8513 & 0.5255 \\
Restaurant (ENG) & 1.2324 & 0.8839 & 0.6094 \\
\hline
\textbf{Average (ENG)} & \textbf{1.2803} & \textbf{0.8676} & \textbf{0.5675} \\
\hline
\end{tabular}
\end{table}

In terms of leaderboard performance, the proposed system achieved a ranking of \textbf{10th position} in the \textbf{laptop domain} and \textbf{11th position} in the \textbf{restaurant domain}. These results demonstrate the competitiveness of the proposed approach in a large-scale shared task setting.

\subsection{Subtask 2 Results}

The proposed approaches for Subtask 2 were evaluated using the official continuous F1 (cF1) metric provided by the SemEval-2026 shared task. The final predictions were submitted to the official CodeBench evaluation platform for both the English laptop and restaurant domains.

\begin{table}[htbp]
\centering
\caption{Official evaluation results for Subtask 2 (DimASTE). Higher cF1 indicates better performance.}
\label{tab:subtask2_results}
\begin{tabular}{lcccc}
\hline
\textbf{Domain} & \textbf{cF1} & \textbf{cPrecision} & \textbf{cRecall} & \textbf{cTP} \\
\hline
Laptop (ENG) & 0.4763 & 0.5215 & 0.4383 & 865.1197 \\
Restaurant (ENG) & 0.5650 & 0.5725 & 0.5577 & 1187.4072 \\
\hline
\textbf{Average (ENG)} & \textbf{0.5206} & \textbf{0.5470} & \textbf{0.4980} & \textbf{1026.2634} \\
\hline
\end{tabular}
\end{table}

To further analyze the effectiveness of the proposed architectures for Subtask 2, a comparative evaluation was conducted between the fully generative \textbf{Llama-only} approach and the proposed \textbf{Hybrid (Llama + RoBERTa)} framework.

\begin{table}[htbp]
\centering
\caption{Comparison between Llama-only and Hybrid approaches for Subtask 2.}
\label{tab:hybrid_comparison}
\begin{tabular}{llccc}
\hline
\textbf{Domain} & \textbf{Method} & \textbf{cPrecision} & \textbf{cRecall} & \textbf{cF1} \\
\hline
\multirow{2}{*}{Laptop}
& Llama-Only & 0.5044 & 0.4439 & \textbf{0.4722} \\
& Hybrid (RoBERTa) & 0.5015 & 0.4412 & 0.4694 \\
\hline
\multirow{2}{*}{Restaurant}
& Llama-Only & 0.6039 & 0.5821 & \textbf{0.5928} \\
& Hybrid (RoBERTa) & 0.5946 & 0.5731 & 0.5836 \\
\hline
\end{tabular}
\end{table}

The comparative results indicate that the fully generative Llama-only approach slightly outperformed the hybrid framework across both domains. In terms of leaderboard rankings, the proposed system achieved \textbf{18th position} in the \textbf{laptop domain} and \textbf{16th position} in the \textbf{restaurant domain}, demonstrating competitive performance in the official SemEval-2026 shared task evaluation.

Although the hybrid approach was designed to leverage the strong regression capabilities of RoBERTa for VA prediction, the experimental results suggest that replacing Llama-generated VA values with RoBERTa predictions did not lead to performance improvements. One possible explanation is that the end-to-end Llama model learns implicit dependencies between aspect terms, opinion expressions, and VA scores during joint generation, whereas separating extraction and regression into independent stages may introduce inconsistencies between components.

Nevertheless, the hybrid framework still demonstrated competitive performance and provided valuable insights into combining generative and discriminative architectures for Dimensional Aspect-Based Sentiment Analysis. These findings highlight the importance of preserving contextual dependencies in structured sentiment generation tasks.

\subsection{Subtask 3 Results}

The proposed baseline approach for Subtask 3 (DimASQP) was evaluated using the official continuous F1 (cF1) metric provided by the SemEval-2026 shared task. The final predictions generated using the \textbf{Llama-3.1-8B-Instruct} baseline model were submitted to the official CodeBench evaluation platform for the English laptop domain.

\begin{table}[htbp]
\centering
\caption{Official evaluation results for Subtask 3 (DimASQP). Higher cF1 indicates better performance.}
\label{tab:subtask3_results}
\begin{tabular}{lcccc}
\hline
\textbf{Domain} & \textbf{cF1} & \textbf{cPrecision} & \textbf{cRecall} & \textbf{cTP} \\
\hline
Laptop (ENG) & 0.3072 & 0.3740 & 0.2606 & 514.6865 \\
\hline
\textbf{Average (ENG)} & \textbf{0.3072} & \textbf{0.3740} & \textbf{0.2606} & \textbf{514.6865} \\
\hline
\end{tabular}
\end{table}


To evaluate the impact of prompt engineering and hyperparameter optimization on quadruplet extraction performance, all three experiments conducted for Subtask 3 were compared using the official evaluation metrics.

\begin{table}[htbp]
\centering
\caption{Comparison of different experiments for Subtask 3 (Laptop Domain).}
\label{tab:subtask3_experiments}
\begin{tabular}{lcccc}
\hline
\textbf{Experiment} & \textbf{Actual F1} & \textbf{cF1} & \textbf{cPrecision} & \textbf{cRecall} \\
\hline
Exp 1: Baseline & 0.3131 & 0.2917 & 0.3511 & 0.2496 \\
Exp 2: Prompt Change & 0.2951 & 0.2749 & 0.3372 & 0.2321 \\
Exp 3: Hyperparameter Optimization & \textbf{0.3302} & \textbf{0.3076} & \textbf{0.3664} & \textbf{0.2651} \\
\hline
\end{tabular}
\end{table}

The comparative results demonstrate that \textbf{Experiment 3}, which focused on hyperparameter optimization, achieved the best overall performance across all evaluation metrics. Furthermore, the final submission based on Experiment 3 achieved \textbf{7th rank} in the official SemEval-2026 Subtask 3 leaderboard for the laptop domain, demonstrating the effectiveness and competitiveness of the proposed optimization strategy in a large-scale shared task setting.

The baseline experiment established a strong initial performance by using a detailed instruction-based prompt with reasoning steps and few-shot examples. However, the second experiment, which simplified the prompt structure, resulted in a performance decrease. This suggests that reducing contextual guidance and example-based instructions may negatively affect the model’s ability to generate accurate and consistent quadruplets.

In contrast, hyperparameter optimization in Experiment 3 significantly improved model performance. Adjustments to training configurations such as learning rate, batch size, sequence length, LoRA parameters, and training epochs contributed to better extraction consistency and improved continuous sentiment prediction.

These findings indicate that, for complex structured generation tasks such as DimASQP, training configuration and optimization strategies may have a greater impact on performance than prompt simplification alone. Furthermore, the experiments highlight the importance of balancing prompt design, model capacity, and fine-tuning strategies when working with instruction-tuned Large Language Models for dimensional sentiment analysis.


\section{Conclusion}
\label{sec:conclusion}

This thesis explored the problem of Dimensional Aspect-Based Sentiment Analysis (DimABSA) within the framework of the SemEval-2026 shared task. Unlike traditional sentiment analysis approaches that rely on discrete sentiment categories, DimABSA aims to capture fine-grained emotional information by combining structured sentiment extraction with continuous valence–arousal (VA) prediction.

Three subtasks were investigated in this work, including dimensional sentiment regression (Subtask 1), triplet extraction (Subtask 2), and quadruplet extraction (Subtask 3). Different transformer-based and Large Language Model (LLM)-based approaches were explored to address the challenges associated with each task.

For Subtask 1, multiple pretrained transformer encoder models were evaluated for continuous sentiment regression. Experimental results demonstrated that \textbf{RoBERTa-large} consistently achieved the best performance across both restaurant and laptop domains. The final system achieved competitive leaderboard rankings, including \textbf{10th place} in the laptop domain and \textbf{11th place} in the restaurant domain, highlighting the effectiveness of transformer-based regression models for dimensional sentiment prediction.

For Subtask 2, two different approaches were explored: a fully generative Llama-based framework and a hybrid architecture combining Llama for textual extraction with RoBERTa for VA regression. The experiments demonstrated that the end-to-end generative approach slightly outperformed the hybrid framework, suggesting that preserving contextual dependencies between extraction and sentiment prediction is crucial for structured sentiment generation tasks. The proposed system achieved \textbf{18th rank} for the laptop domain and \textbf{16th rank} for the restaurant domain in the official shared task evaluation.

For Subtask 3, instruction-tuned LLMs were employed for end-to-end quadruplet extraction. Multiple experiments involving prompt engineering and hyperparameter optimization were conducted. The results showed that hyperparameter optimization had a significant impact on performance and ultimately produced the best system configuration. The final submission achieved \textbf{7th rank} on the official leaderboard for the laptop domain, demonstrating the strong potential of LLM-based approaches for complex structured sentiment extraction tasks.

Overall, the findings of this thesis demonstrate that transformer-based architectures and instruction-tuned LLMs provide effective solutions for DimABSA tasks. The experiments further highlight the importance of prompt design, parameter-efficient fine-tuning, contextual dependency modeling, and hyperparameter optimization in improving structured sentiment extraction performance.

Despite the promising results, several challenges remain open for future research. In particular, improving continuous VA prediction consistency, reducing hallucinated outputs in generative models, and extending the approaches to multilingual and cross-domain settings remain important research directions. Future work may also explore advanced retrieval-augmented generation methods, larger instruction-tuned models, and reinforcement learning-based optimization strategies for dimensional sentiment analysis.

\addcontentsline{toc}{section}{References}
\bibliographystyle{plainurl} 
\bibliography{references}
\end{document}